\section{First-entry ambiguity in the periodic regime}
\label{sec:open-system}

The collision law is the main result of the paper.  This section records
a second consequence of the same cyclic matrix limit: finite observation
of a first-entry residue event.  We state all residue hypotheses
explicitly because the exact path sum and the uniform comparison with
survivor mass require different amounts of nondegeneracy.

\subsection{Bayes error and attainable residues}

Let
\[
  \mathcal F_m:=\sigma(y_0,\ldots,y_{m-1})
\]
be the symbolic prefix filtration.  For an event
\(\mathsf P\in\mathcal F_\infty\), define its depth-\(m\) Bayes error by
\[
  \varepsilon_m(\mathsf P;\nu_\varphi)
  :=\inf\{\nu_\varphi(g\neq\mathbf1_{\mathsf P}):
      g\text{ is }\mathcal F_m\text{-measurable and }g\in\{0,1\}\}.
\]
On the partition into length-\(m\) cylinders, the optimal decision is
the majority decision.  Therefore
\begin{equation}
  \varepsilon_m(\mathsf P;\nu_\varphi)
  =\sum_{a\in A^m}
   \min\{\nu_\varphi(\mathsf P\cap[a]),
          \nu_\varphi([a]\setminus\mathsf P)\}.
  \label{eq:bayes-cylinder}
\end{equation}

Fix an integer \(q\geq2\) and a nonempty proper subset
\[
  \varnothing\neq R\subsetneq\mathbb Z/q\mathbb Z.
\]
The first-entry residue event is
\[
  \mathsf P_R
  :=\{y:\tau_H(y)<\infty,\ \tau_H(y)\bmod q\in R\}.
\]
For \(i\in\mathcal S\), let \(\mathbb P_i\) be the Markov law
conditioned to start at \(i\), and define the positive remaining hit
time
\[
  \tau_H^+:=\inf\{n\geq1:y_n\in A_{\rm hole}\}.
\]
The ambient chain is irreducible and the hole is nonempty, so
\(\mathbb P_i(\tau_H^+<\infty)=1\).  Define
\[
  D_i:=\{n\bmod q:\mathbb P_i(\tau_H^+=n)>0\}
  \subseteq\mathbb Z/q\mathbb Z
\]
and, for \(r\in\mathbb Z/q\mathbb Z\),
\[
  p_{i,r}:=\mathbb P_i(r+\tau_H^+\in R\bmod q),
  \qquad
  b_r(i):=\min\{p_{i,r},1-p_{i,r}\}.
\]

The exact matrix identity below is valid for every \(q\geq2\) and every
nonempty proper \(R\), without a positivity assumption on \(b_r\).  For
the two-sided finite-depth comparison we impose
\begin{equation}
  D_i\cap(R-r)\neq\varnothing,
  \qquad
  D_i\cap((\mathbb Z/q\mathbb Z\setminus R)-r)\neq\varnothing
  \label{eq:residue-nondegeneracy}
\end{equation}
for every \(i\in\mathcal S\) and every
\(r\in\mathbb Z/q\mathbb Z\).  This is equivalent to
\(0<b_r(i)\leq1/2\) for all state-phase pairs.  Finiteness then gives
\[
  \theta:=\min_{i,r}b_r(i)>0.
\]

\subsection{Exact path sum and cyclic amplitudes}

\begin{theorem}[Periodic first-entry Bayes law]
\label{thm:periodic-bayes}
Assume the periodic survivor hypothesis of
Section~\ref{sec:preliminaries}.  Fix \(q\geq2\) and a nonempty proper
\(R\subsetneq\mathbb Z/q\mathbb Z\).  Then, without any residue
nondegeneracy assumption, for every \(m\geq1\),
\begin{equation}
  \varepsilon_m(\mathsf P_R;\nu_\varphi)
  =u_1^{\mathsf T}B_1^{m-1}b_{m-1\bmod q}.
  \label{eq:periodic-bayes-path}
\end{equation}

If in addition \eqref{eq:residue-nondegeneracy} holds for every safe
state and every phase, then for every \(m\geq1\),
\begin{equation}
  \theta\,\nu_\varphi(\tau_H\geq m)
  \leq\varepsilon_m(\mathsf P_R;\nu_\varphi)
  \leq\frac12\nu_\varphi(\tau_H\geq m),
  \label{eq:periodic-bayes-sandwich}
\end{equation}
and
\begin{equation}
  \lim_{m\to\infty}-\frac1m
  \log\varepsilon_m(\mathsf P_R;\nu_\varphi)
  =P_Y(\varphi)-P_{Y_H}(\varphi).
  \label{eq:periodic-bayes-escape}
\end{equation}

Let \(L:=\operatorname{lcm}(d,q)\).  For every
\(t\in\mathbb Z/L\mathbb Z\), as \(m\to\infty\) through
\(m-1\equiv t\pmod L\),
\begin{equation}
  \frac{\varepsilon_m(\mathsf P_R;\nu_\varphi)}
       {\nu_\varphi(\tau_H\geq m)}
  \longrightarrow
  \frac{A_{1,t\bmod d}(b_{t\bmod q})}
       {A_{1,t\bmod d}(\mathbf1)}.
  \label{eq:periodic-bayes-amplitude}
\end{equation}
Every limit in \eqref{eq:periodic-bayes-amplitude} lies in
\([\theta,1/2]\).
\end{theorem}

\begin{proof}
Consider a length-\(m\) cylinder \([a_0\cdots a_{m-1}]\).  If one of
its symbols lies in \(A_{\rm hole}\), then \(\tau_H\leq m-1\) is known
from the observed prefix, so the cylinder contributes zero to
\eqref{eq:bayes-cylinder}.  Suppose instead that all its symbols are
safe.  By the Markov property, the conditional future depends on the
prefix only through the terminal state \(a_{m-1}\).  The current time
phase is \(m-1\bmod q\), so the conditional Bayes ambiguity on this
cylinder is
\[
  b_{m-1\bmod q}(a_{m-1}).
\]
Summing the Markov cylinder weights over all safe words gives
\[
  \varepsilon_m(\mathsf P_R;\nu_\varphi)
  =u_1^{\mathsf T}B_1^{m-1}b_{m-1\bmod q},
\]
proving \eqref{eq:periodic-bayes-path}.  The same path sum with terminal
weight \(\mathbf1\) is
\[
  \nu_\varphi(\tau_H\geq m)
  =u_1^{\mathsf T}B_1^{m-1}\mathbf1.
\]
Under \eqref{eq:residue-nondegeneracy}, the coordinatewise bounds
\(\theta\mathbf1\leq b_r\leq\tfrac12\mathbf1\) prove
\eqref{eq:periodic-bayes-sandwich}.  Formula
\eqref{eq:periodic-bayes-escape} follows from
\eqref{eq:rho-pressure} with \(s=1\).

Finally, if \(m-1\equiv t\pmod L\), then both terminal vectors and
cyclic depth classes are fixed:
\[
  b_{m-1\bmod q}=b_{t\bmod q},
  \qquad
  m-1\equiv t\pmod d.
\]
Apply Proposition~\ref{prop:cyclic-perron} to this terminal vector and
to \(\mathbf1\), and divide the two asymptotics.  Positivity of the
denominator gives \eqref{eq:periodic-bayes-amplitude}; the interval bound
follows either from the coordinatewise bounds or from
\eqref{eq:periodic-bayes-sandwich}.
\end{proof}

\begin{proof}[Proof of Theorem~\ref{thm:intro-bayes}]
This is Theorem~\ref{thm:periodic-bayes} under the explicitly stated
nondegeneracy hypothesis.
\end{proof}

\begin{remark}[What periodicity changes]
When \(d=1\), the cyclic coefficient factors as
\[
  A_{1,0}(v)
  =(u_1^{\mathsf T}r_1)(\ell_1^{\mathsf T}v),
\]
so the common initial factor cancels in
\eqref{eq:periodic-bayes-amplitude}.  For \(d>1\), the initial mass in
each cyclic class is paired with a different terminal class.  The limit
therefore depends on both \(t\bmod d\) and \(t\bmod q\), and a
phase-free quasistationary amplitude need not exist.
\end{remark}

\begin{remark}[Role of the all-state condition]
The exact identity \eqref{eq:periodic-bayes-path} does not use
\eqref{eq:residue-nondegeneracy}.  If some coordinates of \(b_r\)
vanish, Proposition~\ref{prop:cyclic-perron} still gives a positive
leading amplitude whenever
\(A_{1,r_d}(b_{r_q})>0\).  The stronger all-state, all-phase condition is
retained here only for the uniform sandwich
\eqref{eq:periodic-bayes-sandwich}; the statement does not silently use
it for conclusions that require only a leading cyclic coefficient.
\end{remark}
